\Cref{fig:1-challenges}把三个挑战连接到机制与证据。

首先，直接GPU移植可能忽略primary码本的数学结构。
每8-bit子空间索引256个距离，其Cartesian坐标允许精确16+16分解，
但进一步拆分会增加lookup和算术。挑战是利用可分性，同时不假定最小表最快。
候选并行word访问与残差lane合作也需要不同映射。

第二，primary优胜者可能集中在路由全集任何位置。并行组不能安全地各保留
全局需求的一小部分。精确局部保留、全局合并与共享完整分数界提供可靠路线，
却增加缓冲、staging和同步。大batch可能已有足够block，不均衡工作负载仍可能
受益于分组，因而正确性与执行决策耦合。IVF列表选择又是独立于向量保留的原语。

第三，固定倍率可能反复精排无关向量，小倍率则可能过早删除真近邻。
大预算参考一致性无需GT，却不能修复参考自身错误或保证未来查询recall。
样本与预算选择也耗时，必须联合评估质量、校准成本及复用。
这些问题组织迭代项目记录，不把观察收益包装成预注册预测。

\par\medskip\noindent\begin{minipage}{\linewidth}
\centering\captionsetup{hypcap=false}
\resizebox{\linewidth}{!}{\begin{tikzpicture}[
 font=\small\linespread{1}\selectfont,
 box/.style={draw=ThesisNavy!65,rounded corners=2pt,align=center,
 text width=3.5cm,minimum height=1.65cm,inner sep=4pt},
 arr/.style={-{Latex[length=2mm]},draw=ThesisNavy!80,thick}]
\node[font=\bfseries] at (0,1.25) {Challenge};
\node[font=\bfseries] at (4.6,1.25) {Mechanism};
\node[font=\bfseries] at (9.2,1.25) {Evidence};
\node[box] (a) at (0,0) {Primary codes\\Broad scan cost};
\node[box,fill=ThesisNavy!7] (b) at (4.6,0) {Cartesian tables\\Packed word scan};
\node[box] (c) at (9.2,0) {Table / load tests\\Score checks};
\node[box] (d) at (0,-2.2) {Uneven query work\\Global top-$c$};
\node[box,fill=ThesisNavy!7] (e) at (4.6,-2.2) {Grouped scan\\Shared bound\\Global merge};
\node[box] (f) at (9.2,-2.2) {Output checks\\Batch / skew tests};
\node[box] (g) at (0,-4.4) {Refinement effort\\Quality / reuse cost};
\node[box,fill=ThesisNavy!7] (h) at (4.6,-4.4) {Sampled agreement\\Workload budget};
\node[box] (i) at (9.2,-4.4) {Held-out quality\\Calibration cost};
\draw[arr] (a)--(b);\draw[arr] (b)--(c);
\draw[arr] (d)--(e);\draw[arr] (e)--(f);
\draw[arr] (g)--(h);\draw[arr] (h)--(i);
\end{tikzpicture}
}
\captionof{figure}[问题、机制与证据]{问题、机制与证据。箭头组织研究问题，不意味着独立测量的提速比
可以相乘，也不意味着无标签一致性保证GT recall。}
\label{fig:1-challenges}
\end{minipage}\par\medskip
